%% ===================================================================
%% Appendix A. Notation and Preliminaries
%% Central symbol table, standing assumptions (A1)-(A4), the objects the
%% theory reasons about, and the matrix-free evaluation of the operator.
%% Everything downstream references this section instead of redefining.
%% ===================================================================
\section{Notation and Preliminaries}
\label{app:prelim}

This appendix proves every result stated in the main paper and reports the supporting experiments in
full. Its purpose is to let a reader reconstruct each claim end to end, so the proofs are written as
explicit derivations rather than sketches. To keep them readable we fix all notation once here: every
symbol used later appears in \cref{tab:notation}, every abbreviation in the glossary below, and the
standing assumptions and operator identities that recur across proofs are stated a single time and then
invoked by name.

\subsection{Symbols}
\label{app:symbols}

The framework studies one object, the equilibrium of a graph operator, and how it moves when the graph
is perturbed. \Cref{tab:notation} lists the symbols in the order a reader meets them: first the operator
and its two Jacobians, then the sensitivity matrices read off them, then the two perturbation budgets
and the constants that relate them, and finally the classification head and the conformal quantities.

\begin{table*}[t]
\centering
\footnotesize
\setlength{\tabcolsep}{6pt}
\caption{Notation used throughout the appendix. ``First used'' points to the result that introduces or
relies on the symbol.}
\label{tab:notation}
\begin{tabular}{@{}lll@{}}
\toprule
Symbol & Meaning & First used \\
\midrule
\multicolumn{3}{@{}l}{\emph{Operator and Jacobians}}\\
$\Ahat$ & symmetric normalized adjacency $D^{-1/2}(A{+}I)D^{-1/2}$ & \cref{thm:phase_transition}\\
$W$ & tied weight matrix of the layer & \cref{thm:phase_transition}\\
$\phi,\ \diag(\phi')$ & $1$-Lipschitz activation; its derivative mask & (A1)\\
$\zstar$ & equilibrium (fixed point) of $F_\theta$ & \cref{thm:phase_transition}\\
$J_z$ & state Jacobian $\diag(\phi')(\Ahat\otimes W)$ & \cref{thm:phase_transition}\\
$J_A$ & structural Jacobian $\partial F/\partial\mathrm{vec}(A)$ & \cref{thm:phase_transition}\\
$\kappa$ & contraction factor $\norm{J_z}_2<1$ & (A3)\\
\midrule
\multicolumn{3}{@{}l}{\emph{Sensitivity read off the operator}}\\
$S$ & structural sensitivity $(I{-}J_z)^{-1}J_A$ & \cref{thm:phase_transition}\\
$S_c$ & edge-constrained sensitivity $ (I{-}J_z)^{-1}J_A P_c$ & \cref{sec:framework}\\
$P_c,\ P_E$ & projection onto symmetric, edge-supported perturbations & \cref{thm:cf2s}\\
$S_v,\ S_{c,v}$ & block rows of $S,\ S_c$ at node $v$ & \cref{prop:radius}\\
$\sigma_1(\cdot)$ & leading singular value & \cref{prop:attack}\\
$w_k,\ v_k$ & edge weight $[\Ahat]_{ij}$; edge score $\norm{[S_c]_{:,k}}_2$ & \cref{prop:transfer}\\
$d_k$ & true damage of removing edge $k$ & \cref{prop:transfer}\\
\midrule
\multicolumn{3}{@{}l}{\emph{Budgets and the constants that relate them}}\\
$\eglob$ & mask-free safe radius $\max(0,1/\norm{W}_2{-}\norm{\Ahat}_2)$ & \cref{thm:phase_transition}\\
$\ecrit$ & norm certificate $(1{-}\kappa)/\norm{W}_2$ & \cref{thm:phase_transition}\\
$\espec$ & all-active spectral break $1/\rho(W){-}\rho(\Ahat)$ & \cref{thm:cf2s}\\
$\ebreak,\ \ebreakall$ & deployed / all-active contraction-break budget & \cref{thm:cf2s}\\
$\ereach$ & measured nonlinear break budget & \cref{rem:obs_o1}\\
$\gW$ & nonnormality of $W$, $\norm{W}_2/\rho(W)\ge1$ & \cref{thm:cf2s}\\
$\eta$ & Neumann-to-spectral slack (pseudospectral index) & \cref{obs:eta_bound}\\
$\beta=\sigma_E$ & edge-supported mass of the Perron mode $u_1$ & \cref{thm:cf2s}\\
$C$ & bracket constant $\gW(1{+}\kappa)/(1{-}\kappa)$ & \cref{thm:cf2s}\\
$L_J$ & equilibrium-map curvature, $\le\norm{W}_2^2\norm{\zstar}$ & \cref{prop:transfer}\\
\midrule
\multicolumn{3}{@{}l}{\emph{Classification head and conformal prediction}}\\
$f=W\zstar_v{+}b$ & node logits read by the linear head & \cref{prop:radius}\\
$m_v^{(c)}$ & logit margin $f_{y_v}{-}f_c$ to competitor $c$ & \cref{prop:radius}\\
$r_v$ & per-node first-order sensitivity radius & \cref{prop:radius}\\
$\score^{\mathrm{TPS}},\ \score^{\mathrm{APS}}$ & the two conformity scores & \cref{def:conf-scores}\\
$\Cset(v),\ \hat q$ & prediction set; conformal threshold & \cref{def:conf-scores}\\
$L_1^{(c)},\ C_v$ & first-order and curvature score-shift constants & \cref{lem:score-shift}\\
$\alpha,\ n_{\mathrm{cal}}$ & target miscoverage; calibration-set size & \cref{thm:robust-cov}\\
\bottomrule
\end{tabular}
\end{table*}

\noindent\textbf{Abbreviations.}
IFT, implicit function theorem; JVP/VJP, Jacobian--vector and vector--Jacobian products (matrix-free
directional derivatives); SVD/rSVD, (randomized) singular value decomposition; PGD, projected gradient
descent; BCD, block coordinate descent; MC, Monte Carlo; TPS/APS, threshold and adaptive prediction-set
conformity scores. Norms are $\norm{\cdot}_2$ (spectral, on operators; Euclidean, on vectors) and
$\norm{\cdot}_F$ (Frobenius); $\rho(\cdot)$ is the spectral radius and $\otimes$ the Kronecker product.

\subsection{Standing assumptions}
\label{app:assumptions}

All proofs operate on the same operator and in the same certified regime, recorded once here.

\begin{definition}[Operator and certified regime]
\label{ass:tight-v2}
Let $F_\theta(Z,\Ahat)=\phi(\Ahat Z W^\top + X_{\mathrm{proj}})$ with equilibrium
$\zstar=F_\theta(\zstar,\Ahat)$. We assume:
\begin{enumerate}
\item[(A1)] \emph{Lipschitz activation.} $\phi$ is $1$-Lipschitz (ReLU or similar), so
$\norm{\diag(\phi')}_2\le1$ pointwise; the nonsmoothness of ReLU at $0$ is handled by the conservative
calculus of \citet{bolte2021conservative}, which supplies an IFT on each linear region.
\item[(A2)] \emph{Bounded weight.} $\norm{W}_2\le c$ through spectral normalization. This caps the
\emph{numerator} of $\gW$ but does not by itself bound $\rho(W)$ away from $0$ (see
\cref{app:bracket}), so (A2) does not imply (A3).
\item[(A3)] \emph{Trained contractivity.} $\kappa:=\norm{J_z}_2<1$, equivalently
$\norm{\Ahat}_2\norm{W}_2<1$, verified post-training ($\kappa\in[0.14,0.59]$ across our suite).
\item[(A4)] \emph{All-active operating point.} On the operating set every unit is active,
$\phi'\equiv1$, so $J_z=\Ahat\otimes W$. Only the \emph{upper} side of the boundary theorem
(\cref{app:bracket}) invokes (A4); every certificate uses (A1)--(A3) alone.
\end{enumerate}
Unless stated otherwise we work in the certified regime $\varepsilon<\ecrit$, where the perturbed
operator is still a contraction and the resolvent $(I-J_z)^{-1}$ is well defined.
\end{definition}

\subsection{The objects the theory reasons about}
\label{app:objects}

The operator $F_\theta$ maps a node-state matrix $Z$ to a new one by one round of graph propagation and a
weight $W$; its fixed point $\zstar$ is the equilibrium an implicit GNN settles into. Two Jacobians
describe how that equilibrium responds to a perturbation. The \emph{state} Jacobian
$J_z=\diag(\phi')(\Ahat\otimes W)$ measures how the next state depends on the current one; under (A3) it
is a contraction, so a perturbation injected into the loop decays geometrically. The \emph{structural}
Jacobian $J_A=\partial F/\partial\mathrm{vec}(A)$ measures how the operator depends directly on the
graph. Differentiating the fixed-point equation $\zstar=F_\theta(\zstar,\Ahat)$ through the IFT combines
the two into the structural sensitivity
\begin{equation}
S=(I-J_z)^{-1}J_A,
\label{eq:S-def}
\end{equation}
the single object every diagnostic is read from. The inverse resolvent $(I-J_z)^{-1}$ is what turns a
local Jacobian into the equilibrium response: it sums the effect of the perturbation as it circulates
around the loop. Restricting the input to feasible (symmetric, edge-supported) perturbations through the
projection $P_c$ gives the constrained operator $S_c=(I-J_z)^{-1}J_A P_c$ used in practice.

From $S$ and $S_c$ the paper reads three diagnostics, each proved in a later section: the certified
budget $\ecrit$ below which the prediction cannot break (\cref{app:proofs,app:bracket}); the per-node
radius $r_v$ that turns the margin of node $v$ into a safe perturbation distance (\cref{app:proofs}); and
the per-edge score $v_k=\norm{[S_c]_{:,k}}_2$ that ranks how much each edge matters
(\cref{app:transfer}). The conformal certificate of \cref{app:conformal} upgrades the first-order radius
into a distribution-free coverage guarantee.

\subsection{Matrix-free evaluation of the sensitivity operator}
\label{app:matrixfree}

None of the diagnostics forms $S$ explicitly, which is what lets the pipeline scale; the routing is
\cref{alg:aegis} and we record here why a single matrix--vector rule suffices. Applying $S_c$ to a
vector $v$ expands the resolvent as a Neumann series,
\begin{equation}
S_c\,v=(I-J_z)^{-1}J_A P_c\,v=\sum_{j=0}^{\infty}J_z^{\,j}\,J_A P_c\,v,
\label{eq:neumann}
\end{equation}
which converges because $\kappa<1$ (A3). Truncating at depth $K$ leaves a tail bounded by the same
contraction factor,
\begin{equation}
\Big\lVert S_c v-\sum_{j=0}^{K}J_z^{\,j}J_A P_c v\Big\rVert_2
\le\frac{\kappa^{\,K+1}}{1-\kappa}\,\norm{J_A P_c v}_2,
\label{eq:neumann-tail}
\end{equation}
so the depth needed for a target accuracy grows only logarithmically in $1/\kappa$ ($K\in[20,50]$ for
$\kappa<0.8$). Each term $J_z^{\,j}(\cdot)$ is one forward-mode JVP costing $O(Nd)$ memory rather than
the $O((Nd)^2)$ of a dense $S$, and a randomized SVD returns the leading pair $(\sigma_1,v_1)$ that the
attack and defense use. The well-separated singular spectrum (gap $0.39$--$0.50$) is why one rSVD
matches an iterative attacker; the empirical truncation residual $\kappa^{200}\in[10^{-105},10^{-48}]$
on the suite confirms \eqref{eq:neumann-tail} is slack at the depths we run.
